\section*{Availability and implementation}

PyFgsea is released under the MIT License. Source code, documentation, and tutorials are available at \url{https://github.com/shayuanxukuang/pyfgsea}. The Python package can be installed directly from PyPI (\texttt{pip install pyfgsea}) and supports Linux, Windows, and macOS. \rev{A permanent archival snapshot is also available at Zenodo: \href{https://doi.org/10.5281/zenodo.19446446}{10.5281/zenodo.19446446}.}

\section*{Data availability}

\rev{All scripts and minimal example data required to reproduce the reported software benchmarks and rolling-window workflow are provided in the public PyFgsea repository under \texttt{repro/}, \texttt{examples/}, and the package source tree (see Availability and implementation).} Public reference resources used in the showcase analyses include Human Cell Atlas datasets (accessed via HCAData) and MSigDB Hallmark gene sets; access and licensing follow the respective providers' terms.

\section*{Author contributions}

K.W. implemented the software and performed experiments. H.S. supervised the project and contributed to manuscript editing. All authors approved the final manuscript.

\section*{Conflict of interest}

None declared.

\section*{Funding}

None declared.

\begin{thebibliography}{10}

\bibitem{ref1}
Aravind Subramanian, Pablo Tamayo, Vamsi K. Mootha, Sayan Mukherjee, Benjamin L. Ebert, Michael A. Gillette, Amanda Paulovich, Scott L. Pomeroy, Todd R. Golub, Eric S. Lander, and Jill P. Mesirov.
Gene set enrichment analysis: A knowledge-based approach for interpreting genome-wide expression profiles.
\emph{Proceedings of the National Academy of Sciences of the United States of America}, 102(43):15545--15550, 2005.

\bibitem{ref2}
Zhi Fang, Xianbin Liu, and Grant Peltz.
GSEAPY: gene set enrichment analysis in Python.
\emph{Bioinformatics}, 39(1):btac757, 2023.

\bibitem{ref3}
Alexander Lachmann, Daniel J. B. Clarke, Zhuorui Xie, and Avi Ma'ayan.
blitzgsea: efficient computation of gene set enrichment analysis through gamma distribution approximation.
\emph{Bioinformatics}, 38(8):2356--2357, 2022.

\bibitem{ref4}
Gennady Korotkevich, Vladimir Sukhov, Nikolay Budin, Boris Shpak, Maxim N. Artyomov, and Alexey Sergushichev.
Fast gene set enrichment analysis.
\emph{bioRxiv}, 2021.

\bibitem{ref5}
Laleh Haghverdi, Maren B\"uttner, Fabian A. Wolf, Fabian Buettner, and Fabian J. Theis.
Diffusion pseudotime robustly reconstructs lineage branching.
\emph{Nature Methods}, 13(10):845--848, 2016.

\bibitem{ref6}
F. Alexander Wolf, Philipp Angerer, and Fabian J. Theis.
Scanpy: large-scale single-cell gene expression data analysis.
\emph{Genome Biology}, 19(1):15, 2018.

\bibitem{ref7}
Arthur Liberzon, Colin Birger, Helga Thorvaldsd\'ottir, Mahmoud Ghandi, Jill P. Mesirov, and Pablo Tamayo.
The Molecular Signatures Database Hallmark gene set collection.
\emph{Cell Systems}, 1(6):417--425, 2015.

\bibitem{ref8}
Arthur Liberzon, Aravind Subramanian, Reid Pinchback, Helga Thorvaldsd\'ottir, Pablo Tamayo, and Jill P. Mesirov.
Molecular Signatures Database (MSigDB) 3.0.
\emph{Bioinformatics}, 27(12):1739--1740, 2011.

\bibitem{ref9}
Aviv Regev, Sarah A. Teichmann, Eric S. Lander, Ido Amit, Christophe Benoist, Ewan Birney, Bernd Bodenmiller, Peter Campbell, Piero Carninci, Menna Clatworthy, et al.
The Human Cell Atlas.
\emph{eLife}, 6:e27041, 2017.

\bibitem{ref10}
Federico Marini et al.
HCAData: Accessing the datasets of the Human Cell Atlas in R/Bioconductor.
\emph{Bioconductor}, 2020.

\end{thebibliography}
